% ======================================================================
\section{The mathematics: where every equation comes from}\label{sec:math}
% ======================================================================

The paper contains one displayed equation and a dozen implicit ones. None of
them are new --- that is the point. CONCLAVE's contribution is the \emph{composition},
so every piece should be traceable to established work. This section states each
one, says where it comes from, shows where it lives in the code, and explains
what it does operationally.

\begin{center}
\renewcommand{\arraystretch}{1.25}
\begin{tabular}{@{}>{\raggedright\arraybackslash}p{4.5cm}
                 >{\raggedright\arraybackslash}p{5.6cm}
                 >{\raggedright\arraybackslash}p{4.9cm}@{}}
\toprule
\rowcolor{conclavebg}
{\sffamily\bfseries\color{conclaveblue}Equation} &
{\sffamily\bfseries\color{conclaveblue}Origin} &
{\sffamily\bfseries\color{conclaveblue}In the code} \\
\midrule
Shannon entropy & Shannon 1948 & \file{brain/uncertainty.py} \\
Soft-voting aggregation & Bayesian model averaging & \file{brain/ensemble.py} \\
Aleatoric/epistemic split & Houlsby et al.\ 2011 (BALD) & \file{brain/uncertainty.py} \\
Margin & Scheffer et al.\ 2001 & \file{brain/uncertainty.py} \\
Vote entropy & Dagan \& Engelson 1995 & \file{brain/uncertainty.py} \\
Risk--coverage, AURC & El-Yaniv \& Wiener 2010 & \file{brain/uncertainty.py} \\
TF--IDF & Sp\"arck Jones 1972 & \file{brain/featurize.py} \\
Hashing trick & Weinberger et al.\ 2009 & \file{brain/featurize.py} \\
Kernel-weighted prior & Nadaraya--Watson & \file{memory/episodic.py} \\
Linear opinion pool & Stone 1961 & \file{core/agent.py} \\
A$^\star$ with admissible $h$ & Hart et al.\ 1968 & \file{planning/planner.py} \\
Paired bootstrap & Efron 1979 & \file{experiments/exp3\_sequential.py} \\
\bottomrule
\end{tabular}
\end{center}

% ----------------------------------------------------------------------
\subsection{Entropy, and why it is normalised}

Shannon entropy of a distribution $p$ over $K$ actions:
\begin{equation}
\mathcal{H}[p] \;=\; -\sum_{k=1}^{K} p_k \log p_k
\eqsrc{Shannon 1948}
\end{equation}

\noindent measured in nats. It is $0$ when $p$ puts all mass on one action, and
maximal when $p$ is uniform, where $\mathcal{H}[\text{uniform}] = \log K$.

Because CONCLAVE compares uncertainty across problems with different numbers of
actions --- five tools here, ten humanitarian categories there --- every
quantity is divided by $\log K$:
\begin{equation}
\tilde{\mathcal{H}}[p] \;=\; \frac{\mathcal{H}[p]}{\log K} \;\in\; [0,1].
\end{equation}

\noindent Without this the same gate threshold would mean different things in
different deployments. In \file{conclave/brain/uncertainty.py} this is the
\code{logk} divisor appearing in every return value.

% ----------------------------------------------------------------------
\subsection{Aggregating the committee}

With $M$ members producing posteriors $p_1,\dots,p_M$ and weights $w_m$ summing
to one, the ensemble posterior is the weighted mean:
\begin{equation}
\bar p \;=\; \sum_{m=1}^{M} w_m\, p_m .
\label{eq:agg}
\end{equation}

\noindent This is the empirical form of Bayesian model averaging: treat the
members as samples from a posterior over models and marginalise. With
\code{weighting="uniform"} all $w_m = 1/M$, which is what the paper uses and is
hard to beat when members are comparable in quality.

\begin{keybox}[Why averaging \emph{probabilities} and not votes]
Averaging hard votes throws away exactly the information the gate needs. Two
committees can both vote 3--1 for the same action while one is unanimous in
spirit ($0.9,0.9,0.9,0.4$) and the other barely decided
($0.51,0.52,0.53,0.49$). Equation~\eqref{eq:agg} keeps that difference;
majority voting destroys it.
\end{keybox}

% ----------------------------------------------------------------------
\subsection{The decomposition --- the one equation in the paper}

This is the equation that earns the committee its place. Total predictive
uncertainty splits into two parts with different operational meanings:

\begin{equation}
\underbrace{\mathcal{H}[\bar p]}_{\text{total}}
\;=\;
\underbrace{\frac{1}{M}\sum_{m=1}^{M} \mathcal{H}[p_m]}_{\textbf{aleatoric}}
\;+\;
\underbrace{\left(\mathcal{H}[\bar p] - \frac{1}{M}\sum_{m=1}^{M}\mathcal{H}[p_m]\right)}_{\textbf{epistemic}}
\label{eq:bald}
\end{equation}

\subsubsection{Where it comes from}

The right-hand bracket is not an arbitrary residual --- it is the \textbf{mutual
information} between the prediction $y$ and the model identity $\theta$:
\begin{equation}
\mathcal{I}(y;\theta \mid x)
= \underbrace{\mathcal{H}\big[\mathbb{E}_{\theta}\, p(y\mid x,\theta)\big]}_{\mathcal{H}[\bar p]}
- \underbrace{\mathbb{E}_{\theta}\big[\mathcal{H}[p(y\mid x,\theta)]\big]}_{\frac1M\sum_m \mathcal{H}[p_m]}
\end{equation}
using the standard identity $\mathcal{I}(y;\theta) = \mathcal{H}(y) - \mathcal{H}(y\mid\theta)$
and approximating the expectation over $\theta$ by the finite committee. This
quantity is the acquisition function from \textbf{Bayesian Active Learning by
Disagreement} (BALD), introduced by Houlsby, Husz\'ar, Ghahramani and Lengyel in
2011 for choosing which point to label next. CONCLAVE repurposes it: instead of
asking \emph{which example should I label}, it asks \emph{should I act at all}.

Using a deliberately diverse ensemble as the stand-in for a posterior over
models follows Lakshminarayanan et al.\ (2017), who showed ensembles give
uncertainty estimates competitive with far more expensive Bayesian methods.

\subsubsection{Why the epistemic term cannot go negative}

Entropy is concave. By Jensen's inequality, for any weights $w_m$,
\begin{equation}
\mathcal{H}\Big[\sum_m w_m p_m\Big] \;\ge\; \sum_m w_m\, \mathcal{H}[p_m],
\end{equation}
so $\mathcal{I} \ge 0$ always, with equality exactly when all members agree.
The code clamps at zero anyway (\code{max(H\_mean - mean\_H, 0.0)}) purely to
absorb floating-point noise.

\begin{okbox}[Verified numerically]
Over 4{,}000 randomly generated committees (2--6 members, 5 actions), the
identity in \eqref{eq:bald} held to $4.4\times10^{-16}$ and the epistemic term
was non-negative in every single case. The check is reproducible; the
degenerate cases below are the ones worth memorising.
\end{okbox}

\subsubsection{What the two terms actually mean}

\begin{center}
\renewcommand{\arraystretch}{1.3}
\begin{tabular}{@{}>{\raggedright\arraybackslash}p{4.2cm}rrr>{\raggedright\arraybackslash}p{5.0cm}@{}}
\toprule
\rowcolor{conclavebg}
{\sffamily\bfseries\color{conclaveblue}Committee state} &
{\sffamily\bfseries\color{conclaveblue}conf} &
{\sffamily\bfseries\color{conclaveblue}alea} &
{\sffamily\bfseries\color{conclaveblue}epis} &
{\sffamily\bfseries\color{conclaveblue}Right response} \\
\midrule
All agree, confident & 0.90 & 0.36 & \textbf{0.00} & Act \\
All agree, all unsure & 0.33 & \textbf{1.00} & 0.00 & Abstain --- more data will not help \\
Members disagree & 0.34 & 0.33 & \textbf{0.67} & Gather or escalate --- out of competence \\
\bottomrule
\end{tabular}
\end{center}

\noindent Rows two and three are the entire argument for the decomposition.
Both have confidence around $0.33$. A single model reporting only ``33\%
confident'' cannot distinguish them --- yet one is a genuinely ambiguous
situation where no amount of extra sensing helps, and the other is the model
being out of its depth. \code{UncertaintyGate} routes them to
\code{ABSTAIN} and \code{GATHER}/\code{ESCALATE} respectively.

\begin{warnbox}[The honest caveat, which is also in the paper]
Measured on CONCLAVE's own benchmarks, the decomposition did \emph{not} beat plain
maximum-posterior confidence at either selective prediction or shift detection
(AUROC 0.955 for the epistemic term versus 0.960 for confidence on novel text).
The theory is clean and the degenerate cases above are real, but on these tasks
the extra machinery did not pay for itself. Equation~\eqref{eq:bald} is
reported because it is what the framework computes, not because the measurement
vindicated it.
\end{warnbox}

% ----------------------------------------------------------------------
\subsection{Margin and vote entropy}

Two cheaper signals the gate also consults. The \textbf{margin} is the gap
between the top two actions,
\begin{equation}
\mathrm{margin}(\bar p) \;=\; \bar p_{(1)} - \bar p_{(2)},
\eqsrc{Scheffer et al.\ 2001}
\end{equation}
where $\bar p_{(i)}$ is the $i$-th largest component. It catches the case where
the top probability is respectable but a rival is right behind it --- a
situation entropy alone smooths over.

\textbf{Vote entropy} discards the probabilities and looks only at which action
each member would pick:
\begin{equation}
V_k = \frac{1}{M}\sum_{m=1}^{M} \mathbf{1}\!\left[\arg\max_j (p_m)_j = k\right],
\qquad
\tilde{\mathcal{H}}[V]
\eqsrc{Dagan \& Engelson 1995}
\end{equation}
This is the original query-by-committee disagreement measure (Seung, Opper and
Sompolinsky, 1992). It is robust to a single badly-calibrated member, because
one member's overconfidence cannot inflate a hard vote.

% ----------------------------------------------------------------------
\subsection{Selective prediction: the risk--coverage curve}\label{sec:selective}

An agent that may decline to act is a \emph{selective classifier}: a pair
$(f,g)$ where $f$ predicts and $g \in \{0,1\}$ decides whether to act.
Its two figures of merit are
\begin{equation}
\mathrm{coverage} = \mathbb{E}[g(x)],
\qquad
\mathrm{selective\ risk} = \frac{\mathbb{E}\big[\ell(f(x),y)\,g(x)\big]}{\mathbb{E}[g(x)]},
\eqsrc{El-Yaniv \& Wiener 2010}
\end{equation}
that is, the error rate \emph{among the decisions actually acted upon}. Sweeping
the gate threshold traces a curve; the area under it (AURC) summarises how well
the uncertainty signal ranks decisions. \code{risk\_coverage\_curve} computes it,
and CONCLAVE reports AURC $=0.946$ on the sequential task.

\code{UncertaintyGate.calibrate} inverts this: given a target precision, it
searches a $24\times24$ grid over $(\text{confidence\_min},
\text{epistemic\_max})$ for the thresholds that \textbf{maximise coverage}
subject to precision-on-accepted $\ge$ target. That is the standard
constrained form of the selective-classification problem, and it is why the
gate's thresholds are calibrated rather than guessed.

% ----------------------------------------------------------------------
\subsection{Text features: TF--IDF and the hashing trick}

\subsubsection{TF--IDF, exactly as scikit-learn computes it}

For term $t$ in document $d$, with $n$ documents and $\mathrm{df}(t)$ containing $t$:
\begin{align}
\mathrm{tf}'(t,d) &= 1 + \log \mathrm{tf}(t,d) && \text{(\code{sublinear\_tf=True})}\\
\mathrm{idf}(t)   &= \log\frac{1+n}{1+\mathrm{df}(t)} + 1 && \text{(\code{smooth\_idf=True})}\\
x_{t,d} &= \mathrm{tf}'(t,d)\cdot \mathrm{idf}(t), \qquad
x_d \leftarrow x_d / \lVert x_d\rVert_2 &&
\eqsrc{Sp\"arck Jones 1972}
\end{align}

\begin{okbox}[Verified against scikit-learn]
Implementing the three lines above by hand reproduces
\code{TfidfVectorizer(sublinear\_tf=True)} to $10^{-12}$. The $+1$ terms are
smoothing, so an unseen term never divides by zero; the sublinear \code{tf}
stops a word repeated twenty times counting twenty times as much.
\end{okbox}

\noindent Because rows are $L_2$-normalised, the dot product between a query
vector and a tool-specification vector \emph{is} their cosine similarity ---
which is precisely feature \#1 of \code{ToolRouter}.

\subsubsection{The hashing trick, and what it costs}

Storing a vocabulary makes the model grow with the corpus. Instead, a hash
function $h : \text{term} \to \{1,\dots,D\}$ maps terms directly to indices:
\begin{equation}
\phi_j(x) \;=\; \sum_{i \,:\, h(i)=j} x_i
\eqsrc{Weinberger et al.\ 2009}
\end{equation}

\noindent The feature space is fixed at $D$ regardless of how much text is seen
--- there is no stored vocabulary at all. This is why \code{n\_word\_features}
sets the artefact size directly, and it is the single most important property
for microcontroller-class deployment.

\begin{warnbox}[CONCLAVE uses the unsigned variant, which is biased]
Weinberger et al.'s analysis relies on a \emph{signed} hash, $\xi(i)\in\{\pm1\}$,
which makes collisions cancel in expectation and preserves inner products
without bias. CONCLAVE sets \code{alternate\_sign=False}, so colliding terms always
add \emph{positively} and similarities are biased slightly upward.

This is a deliberate trade: \code{ComplementNB} requires non-negative inputs,
and keeping a generative member in the committee was judged worth the bias.
With $D = 2^{16}$ collisions are rare enough that it does not show up in the
measurements, but it is a real approximation and worth knowing if you shrink
$D$ aggressively.
\end{warnbox}

% ----------------------------------------------------------------------
\subsection{Episodic memory as a kernel-weighted prior}

Given a query key $z$ and retrieved cases $(z_i, a_i, r_i)$ with similarity
$s_i = \max(0, 1 - d_i)$, the memory produces a prior over actions:
\begin{equation}
\pi_k \;\propto\; \sum_{i \,:\, a_i = k,\; r_i > 0} s_i \, r_i,
\qquad
\pi \leftarrow \pi^{1/T} \big/ \textstyle\sum_k \pi_k^{1/T}
\end{equation}

\noindent This is a Nadaraya--Watson kernel-weighted average with a
truncated-cosine kernel, restricted to cases that were actually rewarded, and
$T$ a temperature that sharpens ($T<1$) or flattens ($T>1$) the result. When no
retrieved case has positive reward it returns the uniform distribution, so
memory never biases a decision it knows nothing about.

The result is combined with the ensemble posterior as a \textbf{linear opinion
pool}:
\begin{equation}
p_{\text{final}} \;=\; (1-\omega)\,\bar p \;+\; \omega\,\pi,
\qquad \omega = \code{memory\_weight} = 0.25
\eqsrc{Stone 1961}
\end{equation}

\noindent A linear pool, unlike a logarithmic one, cannot veto: if either source
assigns an action non-negligible mass it survives. For an agent that must not be
talked out of a correct action by a stale memory, that is the safer choice.

% ----------------------------------------------------------------------
\subsection{The planner, and an honest caveat}

\code{Planner.plan} is A$^\star$ over sets of predicates. With $g(s)$ the cost so
far and $h(s)$ an estimate of the cost remaining, it expands the node minimising
$f(s) = g(s) + h(s)$, and A$^\star$ returns an optimal plan \emph{provided}
$h$ never overestimates --- the admissibility condition of Hart, Nilsson and
Raphael (1968).

CONCLAVE uses the goal-count heuristic:
\begin{equation}
h(s) \;=\; \big|\,G \setminus s\,\big|
\end{equation}
the number of goal predicates not yet satisfied.

\begin{dangerbox}[This heuristic is not admissible in general]
The docstring calls it admissible. That holds only under two assumptions:
\begin{enumerate}[leftmargin=*,itemsep=1pt,topsep=2pt]
  \item every action adds \textbf{at most one} goal predicate, and
  \item every action costs \textbf{at least 1}.
\end{enumerate}
Break either and $h$ can overestimate, and A$^\star$ may return a suboptimal
plan. A concrete case, reproducible in three tools:
\begin{lstlisting}[style=py]
step_a     : effects {a},       cost 0.1
finish_bc  : requires a, effects {b,c}, cost 0.5   # total 0.6  <- optimal
direct     : effects {a,b,c},   cost 1.0
# goal = {a,b,c}   ->   CONCLAVE returns ['direct'], cost 1.0
\end{lstlisting}
At state $\{a\}$ we have $g=0.1$ and $h=|\{b,c\}|=2$, so $f=2.1$ --- but the
true remaining cost is $0.5$. The optimal path is pushed behind the
\code{direct} goal node at $f=1.0$, which is popped first.
\end{dangerbox}

\noindent The shipped disaster tools all have single-predicate effects and unit
costs, so the guarantee holds there. If you add a multi-effect or sub-unit-cost
tool, either accept ``good but not provably optimal'', or restore admissibility
with
\begin{equation}
h(s) \;=\; \big|G \setminus s\big| \cdot \frac{c_{\min}}{e_{\max}},
\end{equation}
where $c_{\min}$ is the cheapest action cost and $e_{\max}$ the most goal
predicates any single action can add. Setting $h \equiv 0$ turns A$^\star$ into
uniform-cost search: always optimal, just slower.

% ----------------------------------------------------------------------
\subsection{The paired bootstrap in the headline experiment}

The connectivity result compares two policies on \emph{the same} episode seeds,
so the comparison is paired. For per-episode differences
$d_i = R_i^{\text{CONCLAVE}} - R_i^{\text{cloud}}$, resample indices with
replacement $B=4000$ times and take the empirical $2.5$th and $97.5$th
percentiles of $\bar d^{*}$:
\begin{equation}
\bar d^{*(b)} = \frac{1}{n}\sum_{i=1}^{n} d_{\,\iota^{(b)}_i},
\qquad
\text{CI}_{95} = \Big[\,Q_{2.5}\big(\bar d^{*}\big),\; Q_{97.5}\big(\bar d^{*}\big)\Big]
\eqsrc{Efron 1979}
\end{equation}

\noindent Pairing matters here: episode difficulty varies enormously (when the
surge arrives, how big it is), and pairing removes that variance instead of
letting it swamp the effect. It is also what lets the paper make its most
important negative claim --- that at zero outage the two are
\emph{indistinguishable}, $\Delta = 0.6$ with CI $[-0.7, 1.9]$ --- which is the
check that the comparison was fair in the first place.

% ----------------------------------------------------------------------
\subsection{Reading list, in dependency order}

\begin{enumerate}[leftmargin=*,itemsep=3pt]
  \item \textbf{Shannon (1948)}, \emph{A Mathematical Theory of Communication} --- entropy.
  \item \textbf{Sp\"arck Jones (1972)} --- inverse document frequency.
  \item \textbf{Hart, Nilsson \& Raphael (1968)} --- A$^\star$ and admissibility.
  \item \textbf{Seung, Opper \& Sompolinsky (1992)}; \textbf{Dagan \& Engelson (1995)} --- query by committee, vote entropy.
  \item \textbf{Breiman (2001)} --- random forests; where ensemble diversity comes from.
  \item \textbf{Rennie et al.\ (2003)} --- complement naive Bayes.
  \item \textbf{Weinberger et al.\ (2009)} --- the hashing trick and its guarantees.
  \item \textbf{El-Yaniv \& Wiener (2010)} --- selective classification.
  \item \textbf{Houlsby et al.\ (2011)} --- BALD; equation~\eqref{eq:bald}. \emph{Start here if you read only one.}
  \item \textbf{Lakshminarayanan et al.\ (2017)} --- ensembles as practical uncertainty.
\end{enumerate}
